\newcommand{\single}[1]{$\langle{\mbox{\em #1\/}}\rangle$}
\newcommand{\one}[1]{$\langle{\mbox{\em #1\/}}\rangle^\dagger$}
\newcommand{\zero}[1]{$\langle{\mbox{\em #1\/}}\rangle^\star$}
\newcommand{\reds}{$\stackrel{1}{\longrightarrow}$}
\newcommand{\red}{${\longrightarrow}$}


\chapter{THOR}
\label{chapter-THOR}

	This chapter presents THOR's syntax and
gives an informal semantics to its built-in primitives.

\section{Syntax}

	THOR is a small language with a minimal syntax.  The domain 
of THOR is the set of well-formed lambda expressions,
as defined by the following BNF formulae:
\begin{tabbing}
12345\=\kill
\>	\single{expression} ::= \one{exp}\\[.10in]
\>	\single{exp} ::= \= (\single{expression}) $\mid$ \single{abstraction}
				$\mid$ \single{constant} $\mid$ \single{list} $\mid$ \\
\>\>			 \single{structure} $\mid$ \single{string}\\[.10in]
\>	\single{abstraction} ::= \=
		(LAMBDA (\one{symbol}) \single{expression}) $\mid$\\
\>\>		(LET (\one{binding-pair}) \single{expression}) $\mid$\\
\>\>		(LET* (\one{binding-pair}) \single{expression}) $\mid$\\
\>\>		(LETREC (\one{binding-pair}) \single{expression}) \\[.10in]
\>	\single{binding-pair} ::= (\single{symbol} \single{expression})\\[.10in]
\>	\single{constant} ::= \single{character} $\mid$  \single{integer} $\mid$
		 \single{float} $\mid$ \single{symbol}\\[.10in]
\>	\single{list} ::= [ \zero{exp}]
		 $\mid$ [ \one{exp} $\parallel$ \single{exp}]\\ [.10in]
\>	\single{structure} ::= \{ \single{symbol} \zero{exp} \}\\[.10in]
\>	\single{string} ::= ``\zero{character}''\\[.10in]
\end{tabbing}
A superscript $\dagger$ means one or more occurrences and a 
superscript $\star$ means zero or more occurrences.

	As is customary in functional languages, application is 
left associative and is indicated simply by juxtaposition.
Parentheses can be used to group applications and override left associativity.

	THOR is extremely permissive about the content of expressions it
will reduce.  For the most part, THOR has no real semantic notion of 
`content' or `meaning'. If an expression is syntactically well-formed, 
then it is allowable.  For example, the expression {\mbox {\tt (1 2)}}, 
which is the integer 1 applied to the integer 2, is allowable, and
reduces to itself.  This
expression, which has no meaning in conventional programming languages
(indeed, it would almost surely generate a compilation or run-time
error), is perfectly acceptable in THOR.  You are free to assign to it 
whatever meaning you wish.

	A result of this permissiveness is that there are no run-time errors 
in THOR.  If there is a problem of some kind, then reduction simply does
not take place.  For example, if the arguments to a primitive do not
have the proper type, then the primitive does not fire.  

	This may all seem
crazy at first, but that is only because you are used to using conventional
programming languages which are highly restrictive in their view of what
computation is.  


\section{Abstraction Mechanisms}

		Abstraction is the way functions are formed in the lambda
calculus.  THOR provides four abstraction mechanisms:  {\tt LAMBDA}, 
{\tt LET}, {\tt LET*}, and {\tt LETREC}.  The first three are all simple
variants of lambda binding, and the fourth, {\tt LETREC}, will be
discussed in a later section devoted to recursion.

\subsection{\tt LAMBDA}

\begin{quote}
	{\tt (LAMBDA (${\mbox{\sl var}}_1$ \ldots) {\sl body})}
\end{quote}

	{\tt LAMBDA} is the basic abstraction mechanism. It's semantics are 
equivalent to $\lambda$, but its syntax differs in that a single 
{\tt LAMBDA} can be used to bind more than one variable.  For example, 
the expression 
{\mbox ($\lambda$x.$\lambda$y.(+ x y))} can be written as 
{\mbox{\tt (lambda (x y) (+ x y))}} in THOR.

	When a {\tt LAMBDA} abstraction is applied to an argument, a
$\beta$-reduction can occur.  When this happens, the argument is
substituted for each free occurrence of the bound variable  in
the body of the abstraction.  The lambda variables are  
bound from left to right in sequence, as illustrated in this example:
	\footnote{The $\longrightarrow$ symbol is used to indicate the 
		expression on the left reduces to the expression on the right.  
		When accompanied by an integer, the integer indicates how many 
		reductions have taken place. }
\begin{quote}
\begin{tt}
	((lambda (x y) (+ x y)) 3 7) \reds \\
	((lambda (y) (+ 3 y) 7) \reds \\
	(+ 3 7) \reds \\
	10
\end{tt}
\end{quote}

	It may also be the case that a lambda variable is not part of a 
$\beta$-redex, so it does not get bound.  In the rest of this manual,
such a variable will be referred to as an {\em unbound variable}.


\subsection{\tt LET}

\begin{quote}
	{\tt (LET ((${\mbox{\sl var}}_1$ ${\mbox{\sl exp}}_1$) \ldots) 
				{\sl body})}
\end{quote}

	{\tt LET} is a sugared form of $\lambda$ binding, in which
all the binding pairs are bound in the same context. 
The expression 
\begin{center}
{\mbox{\tt (let (($v_1$ $e_1$) \ldots\ ($v_n$ $e_n$)) exp)}} 
\end{center}
is equivalent to 
\begin{center}
{\mbox{\tt ((lambda ($v_1$ \ldots\ $v_n$) exp) $e_1$ \ldots\ $e_n$)}}
\end{center}
Each binding pair is a $\beta$-redex, and they are contracted from left
to right:
\begin{quote}
\begin{tt}
	(let ((a 1) (b 2)) (+ a b)) \reds \\
	(let ((b 2)) (+ 1 b)) \reds \\
	(+ 1 2) \reds \\
	3
\end{tt}
\end{quote}

	Although {\tt LET} can be expressed in terms of {\tt LAMBDA}, it is
included in THOR because in many circumstances it is more
understandable than an equivalent {\tt LAMBDA} expression.


\subsection{\tt LET*}

\begin{quote}
	{\tt (LET* ((${\mbox{\sl var}}_1$ ${\mbox{\sl exp}}_1$) \ldots) 
				{\sl body})}
\end{quote}

	{\tt LET*} is a sugared form of $\lambda$ binding, where
the binding pairs are bound in nested contexts.
The expression 
\begin{center}
{\mbox{\tt (let* (($v_1$ $e_1$) \ldots\ ($v_n$ $e_n$)) exp)}} 
\end{center}
is equivalent to 
\begin{center}
{\mbox{\tt ((lambda ($v_1$) (\ldots\ ((lambda ($v_n$) exp) $e_n$) 
										\ldots)) $e_1$)}}
\end{center}
Each binding pair is a $\beta$-redex, and they are contracted from left
to right:
\begin{quote}
\begin{tt}
	(let ((a 1) (b (+ a 4)) (+ a b)) \reds \\
	(let ((b (+ 1 4)) (+ 1 b)) \reds \\
	(+ 1 (+ 1 4)) $\stackrel{2}{\longrightarrow}$\\
	6
\end{tt}
\end{quote}

	 Although {\tt LET*} can be expressed in 
terms of {\tt LAMBDA}, it is
included in THOR because it is often more understandable.



\section{Constants}

	There are four types of constant in THOR: characters, 
integers, floating point
numbers, and symbols.  With the exception of symbols that
are associated with expressions (see Section~\ref{sec-defining}),
constants reduce to themselves.

	All constants are represented by character sequences that are
delimited by the following: white space (spaces, tabs,
and newlines), fence characters (parens, braces, and brackets), commas,
semicolons, single quotes, double quotes, and the vertical bar.


\subsection{Characters}

	Character constants are represented using the C language
conventions --- the character is enclosed in single quotes, and the
backslash is used as an escape character for delimiters.  The escaped
characters are:
\begin{tabbing}
12345\=12345\=\kill
\>\verb1\n1 \>Newline\\
\>\verb1\t1 \>Tab\\
\>\verb1\b1 \>Backspace\\
\>\verb1\r1 \>Carriage Return\\
\>\verb1\f1 \>Form Feed\\
\>\verb1\'1 \>Single Quote\\
\>\verb1\"1 \>Double Quote\\
\>\verb1\\1 \>Backslash
\end{tabbing}

Unlike C,
however,	characters are represented as a separate data type in THOR, 
distinct from the other basic types.


\subsection{Character Primitives}

	At the present time, there are four built-in primitives applicable to
characters: 
\begin{tabbing}
{\tt char->int} {\sl c}12345\=\kill
{\tt ATOM?} {\sl c} \> {\tt TRUE} if {\sl c} is a character; does not fire
if {\sl c} \\ \> is an unbound variable.\\
{\tt CHAR->INT} {\sl c} \> converts character {\sl c} to an integer\\ 
{\tt INT->CHAR} {\sl x} \> converts integer {\sl x} to a character\\
{\tt CHAR?} {\sl c}     \> {\tt TRUE} if {\sl c} is a character; does not fire
if {\sl c} \\ \> is an unbound variable; {\tt FALSE} otherwise.
\end{tabbing}
All four are {\em strict} unary operators.  A strict operator is one which 
reduces its arguments before it is fired.  



\subsection{Integers}

	Integers are represented by an optional sign followed by one or more 
decimal digits.  (Currently the only radix supported is base 10.)
The value of the largest and smallest integers available
varies from machine to machine, but will at least be equivalent to that
of a signed long in a machine's C language implementation.

\subsection{Floating Point Numbers}

	Floating point numbers are represented by an optional sign followed
by zero or more decimal digits, followed by a period, then one or more
decimal digits. (Exponential notation is not currently supported.)  The
range and precision of floating point numbers available varies from
machine to machine, but will at least be equivalent to that of a single
precision float in a machine's C language implementation.



\subsection{Numeric Primitives}

	THOR includes most of the common unary and binary arithmetic and
relational operators.  All of the numeric primitives are strict, and
integer $\leftrightarrow$ floating point conversions
are automatically performed if necessary. 
If there are not enough arguments or the arguments
are not of numeric type, the operator is not fired.  The firing of an
operator counts as one reduction.

\subsubsection{Unary Arithmetic Operators}
\begin{tabbing}
{\tt CEILING} {\sl x}12345\=\kill
{\tt 1+} {\sl x} \>	Increments {\sl x} by 1\\
{\tt 1-} {\sl x} \>  Decrements {\sl x} by 1\\
{\tt ABS} {\sl x} \> Absolute value of {\sl x}\\
{\tt CEILING} {\sl x} \>	The smallest integer value greater than or
			equal to {\sl x}\\
{\tt FLOOR} {\sl x} \>	The largest integer value less than or equal to
			{\sl x}\\
{\tt MINUS} {\sl x} \>  Changes the sign of {\sl x}\\
{\tt SQRT} {\sl x} \> Square root of {\sl x}
\end{tabbing}


\subsubsection{Binary Arithmetic Operators}
\begin{tabbing}
{\tt MODULO} {\sl x} {\sl y}1234\=\kill
{\tt +} {\sl x} {\sl y} \> Sum of {\sl x} and {\sl y}\\
{\tt -} {\sl x} {\sl y} \> Difference of {\sl x} and {\sl y}\\
{\tt *} {\sl x} {\sl y} \> Product of {\sl x} and {\sl y}\\
{\tt /} {\sl x} {\sl y} \> Quotient of {\sl x} and {\sl y}\\
{\tt EXPT} {\sl x} {\sl y} \>{\sl X} raised to the power of {\sl y}.
		Neither argument can be zero;\\
		\> when {\sl x} is negative, {\sl y} must be an integer.\\
{\tt MAX} {\sl x} {\sl y} \> The larger of the values of {\sl x} and
{\sl y}\\
{\tt MIN} {\sl x} {\sl y} \> The smaller of the values of {\sl x} and
{\sl y}\\
{\tt MODULO} {\sl x} {\sl y} \> The integer remainder of {\sl x} divided
by {\sl y}\\
\end{tabbing}

\subsubsection{Unary Relational Operators}
\begin{tabbing}
{\tt NEGATIVE?} {\sl x}12345\=\kill
{\tt EVEN?} {\sl x} \> {\tt TRUE} if {\sl x} is an even number; {\tt FALSE}
			otherwise\\
{\tt NEGATIVE?} {\sl x} \> {\tt TRUE} if {\sl x} $<$ 0; {\tt FALSE} otherwise\\
{\tt ODD?} {\sl x} \> {\tt TRUE} if {\sl x} is an odd number; {\tt FALSE}
			otherwise\\ 
{\tt POSITIVE?} {\sl x} \> {\tt TRUE} if {\sl x} $\geq$ 0; {\tt FALSE} otherwise
\end{tabbing}

\subsubsection{Binary Relational Operators}
\begin{tabbing}
{\tt $=$} {\sl x} {\sl y}12345\=\kill
{\tt $=$} {\sl x} {\sl y} \> {\tt TRUE} if {\sl x} equals {\sl y}; {\tt FALSE}
otherwise\\ 
{\tt $<>$} {\sl x} {\sl y} \> {\tt TRUE} if {\sl x} is not equal to {\sl y};
{\tt FALSE} otherwise\\ 
{\tt $>$} {\sl x} {\sl y} \> {\tt TRUE} if {\sl x} is greater than {\sl y};
{\tt FALSE} otherwise\\ 
{\tt $<$} {\sl x} {\sl y} \> {\tt TRUE} if {\sl x} is less than {\sl y};
{\tt FALSE} otherwise\\ 
{\tt $>=$} {\sl x} {\sl y} \> {\tt TRUE} if {\sl x} is greater than or equal
to {\sl y}; {\tt FALSE} otherwise\\ 
{\tt $<=$} {\sl x} {\sl y} \> {\tt TRUE} if {\sl x} is less than or equal
to {\sl y}; {\tt FALSE} otherwise
\end{tabbing}

\subsubsection{Type Predicates}

	There are three type predicates that are germane to numbers: 
\begin{tabbing}
{\tt INTEGER?} {\sl x}12345\=\kill
{\tt ATOM?} {\sl x}  \>  {\tt TRUE} if {\sl x} is an integer or a
floating point number;\\ \> does not fire if {\sl x} is
an unbound variable.\\
{\tt INTEGER?} {\sl x}  \> {\tt TRUE} only if {\sl x} is an integer; does not fire
 if {\sl x}\\ \> is an unbound variable; {\tt FALSE} otherwise.\\
{\tt FLOAT?} {\sl x} \> {\tt TRUE} only if {\sl x} is an floating point
number;\\ \> does not fire if {\sl x} is an unbound variable; 
{\tt FALSE} otherwise.
\end{tabbing}


\subsection{Symbols}

	Symbols are the major component of programs in THOR; they can play the
role of variable, primitive, or constant depending on the context.
Symbols are represented by one or more consecutive non-delimiting characters.
Note that this means symbols can begin with digits (such as the
built-in primitive {\tt 1+}) and punctuation marks.  Symbols are not
case sensitive; the HORSE reader converts them internally to all upper-case.

	There is a small group of reserved symbols in THOR which cannot
be used as constants or variables without protecting them first (see
Section~\ref{sec-protect} for an explanation of `protecting' symbols). 
They are {\tt LAMBDA, LET, LET*, LETREC, DEFINE, RDEFINE,} and 
{\tt DEFSTRUCT}.

\subsection{Symbol Primitives}

	There are two primitives that deal with symbols, and both are strict
unary type predicates: 
\begin{tabbing}
{\tt SYMBOL?} {\sl s}12345\=\kill
{\tt ATOM?} {\sl s} \> {\tt TRUE} if {\sl s} is a symbol; does
not reduce if {\sl s} \\ \> is an unbound variable. \\
{\tt SYMBOL?} {\sl s} \> {\tt TRUE} only if {\sl s} is a symbol; does
not reduce if {\sl s} \\  \>is an unbound variable; {\tt FALSE}
otherwise.
\end{tabbing}



\section{Protecting Symbols and Variables}
\label{sec-protect}

	The scoping rules of the lambda calculus specify that a variable is
bound to the smallest enclosing abstraction with a lambda binding of 
the same name. Unfortunately, this can sometimes cause {\em name capture}
problems, and is the 
reason that 
axiomatizations of the lambda calculus generally include a rule that
$\alpha$-substitution must take place before $\beta$-reduction if the
argument contains free variables with the same name as the lambda variable in
the $\beta$-redex.  

Name capture is illustrated by the following reduction sequence:
\begin{quote}\begin{tt}
	(lambda (a b) a b) (lambda (a b) a b) \reds\\
	(lambda (b) ((lambda (a b) a b) b)) \reds\\
	(lambda (b) (lambda (b) b b))
\end{tt}\end{quote}
This result is not what was intended by the original expression;  a
correct result would be $\alpha$-equivalent to 
{\mbox{\tt (lambda (x) (lambda (y) x y))}}.  

	Instead of introducing new variable names via $\alpha$-substitution,
THOR solves the name capture problem with a protection device (represented as
the character \#) which `blocks' a level of unwanted scope binding. 
Using protection devices, THOR's solution to the previous example would
be 
\begin{center}
	{\mbox{\tt (lambda (b) (lambda (b) \#b b))}}
\end{center}
Each \# blocks one
level of scope binding;  multiple \#'s can be used to block several
levels.  For example, the following expressions are $\alpha$-equivalent:
\begin{tt}
\begin{tabbing}
12345\=12345\=\kill
\>	(lambda (x x x) x \#\#x \#x)  $\stackrel{\alpha}{=}$ 
		(lambda (a b c) c a b)\\
\>	(lambda (x) (x (lambda (x) \#x x))) $\stackrel{\alpha}{=}$ \\
\>\>		(lambda (a) (a (lambda (b) a b)))\\
\>	(lambda (x) x \#x) $\stackrel{\alpha}{=}$ 
		(lambda (a) a x)
\end{tabbing}
\end{tt}
In the last expression above,	the protection device is used to declare
that the \underline{symbolic constant} {\tt x} is the desired meaning.  If more \#'s are
used than there are binding levels, the excess \#'s are ignored.




\section{Compound Data Objects}

	There are three classes of compound data objects in THOR: lists, strings,
and structures.  The first two classes can contain a variable number of elements
which can be accessed only in sequence.  The third class are tuples of
fixed size whose elements can be accessed randomly.  All three 
classes of object are constructed in a {\em lazy} manner, which 
means their elements
are not reduced at the time the object is created, but only when they
are accessed.

	The  current implementation of THOR uses special lambda expressions
to represent all three classes of compound data objects.  The HORSE
reader and printer recognize these lambda expressions and compile or
print them using an alternative syntax.  If you use
these objects in conventional ways (i.e., only as arguments to
primitives) you won't ever notice they are lambda expressions.  However,
if one of them becomes the operator in an expression, it behaves like a
lambda expression, and weird things can happen.  Chances are you will
never have to be concerned about this, but be aware there is more to
these objects than apparent from their syntax.


\subsection{Lists}

	Lists are binary recursive data structures which form linearly
ordered collections of values.  The empty list, which contains no values, is
written as either the symbol {\tt NIL} or as {\tt [ \ ]}.  Nonempty
lists are written as {\tt [ {\sl x} $\parallel$ {\sl y}]}, where {\sl x}
is an element of the list, and {\sl y} is recursively the rest of the
list.  (The $\parallel$ is the vertical bar character on most
computer keyboards.)
For example, the list of the first three natural numbers can be
written as
\begin{center}
	{\tt [1 $\parallel$ [2 $\parallel$ [3 $\parallel$ NIL]]]}
\end{center}
However, there is a much better shorthand notation for proper lists (which end
in {\tt NIL}), which omits the intermediate $\parallel$'s and
brackets.  In this shorthand, the above example would be written as
{\mbox{\tt [1 2 3]}}.


\subsection{List Primitives}

	All of the list primitives except for {\tt PAIR?} are implemented as
special lambda expressions, so they take more than one reduction to
accomplish their task and generate intermediate expressions if they are
not finished when the allowed reductions are exhausted.  Here are the
list primitives, along with the number of reductions they require:
\begin{tabbing}
{\tt CON0S} {\sl a} {\sl b}12345\=\kill
{\tt CAR} {\sl l} \> Selects the first element in list {\sl l}.\\ 
\>Costs 5 reductions.\\
{\tt CDR} {\sl l} \> Selects all but the first element of list {\sl
l}.\\
\> Costs 5 reductions.\\
{\tt CONS} {\sl a} {\sl b} \> Constructs the list   
[{\sl a} $\parallel$ {\sl b}]. \\ \> Costs 3 reductions.\\
\end{tabbing}
\pagebreak
\begin{tabbing}
{\tt CON0S} {\sl a} {\sl b}12345\=\kill
{\tt NULL?} {\sl l} \> {\tt TRUE} only if {\sl l} is the symbol {\tt NIL},
which is the empty list;\\ \> does not fire if {\sl l} is an unbound
variable; {\tt FALSE} otherwise.\\
\>Costs 2 reductions.\\
{\tt PAIR?} {\sl l} \> {\tt TRUE} only if {\sl l} is a list; does not
fire if {\sl l} \\ \> is an unbound variable; {\tt FALSE} otherwise.\\
\>Costs 1 reduction.
\end{tabbing}

\subsection{Strings}

	Strings are actually lists of characters.  Their textual
representation is as a sequence of characters enclosed in double quotes.
The backslash character may be used as an escape, just as it is used with
character constants.


\subsection{String Primitives}

	Since strings are actually lists of characters, all of the list
primitives except for {\tt PAIR?} will work with strings also.  There
are currently only two string specific functions:  
\begin{tabbing}
{\tt CONS-STRING} {\sl c} {\sl s}12345\=\kill
{\tt CONS-STRING} {\sl c} {\sl s} \> Creates a new string by prepending
the character {\sl c} \\ \>to the string {\sl s}.\\
{\tt STRING?} {\sl s} \> {\tt TRUE} only if {\sl s} is a string;
does not fire if {\sl s} \\ \> is an unbound variable; {\tt FALSE} otherwise.
\end{tabbing}


\subsection{Structures}

	Structures are tuples that carry a type tag.  They are of fixed
length, and their elements are filled in when they are created.  
Section~\ref{sec-defstruct} gives information about how to use
the {\tt defstruct} special form to define and access structures.
In addition, you can also
create a structure directly by using its print syntax.  For example, when
\begin{center}
	{\tt \{binop + 5 6\}}
\end{center}
is read by the HORSE reader, a structure of type {\tt binop} is created
with {\tt +} as its first element, {\tt 5} as its second element, and
{\tt 6} as its third element.  Using the print syntax is convenient, but
there is no check to see if a structure created in this way conforms to
the definition of structures of that type (if such a definition has
been made).


\subsection{Structure Primitives}

	In addition to the structure defining, creating, and accessing
mechanisms described in Section~\ref{sec-defstruct}, there are two
primitives for operating on structures:  
\begin{tabbing}
{\tt STRUCTURE?} {\sl s}12345\=\kill
{\tt STRUCTURE?} {\sl s} \> {\tt TRUE} only if {\sl s} is a structure;
does not fire \\ \> if {\sl s} is an unbound variable; {\tt FALSE}
otherwise.\\
{\tt TAG} {\sl s} \> Reduces to the type tag of structure {\sl s}.
\end{tabbing}
	

\section{Logical Primitives}

	The logical operators are {\tt AND, OR,} and {\tt NOT}.  They operate
on boolean values, which are represented by the symbols {\tt TRUE} and
{\tt FALSE}.  

	The {\tt NOT} is a strict unary operator with simple semantics: 
\begin{quote}\begin{tt}
	NOT TRUE \reds\  FALSE\\
	NOT FALSE \reds\  TRUE
\end{tt}\end{quote}
The {\tt AND} and {\tt OR} can take any number of arguments, and their
operational semantics are more complicated.  

\subsection{\tt AND}

	The semantics of the {\tt AND} primitive can be expressed by the
following rules:
\begin{quote}
	{\tt AND TRUE \reds\  TRUE}\\
		 If {\tt AND} has a single argument which
reduces to {\tt TRUE}, then the {\tt AND} expression also reduces to
{\tt TRUE}\\[.10in]
	{\tt AND $e_1$ \ldots\ TRUE \ldots\ $e_n$ \reds\  AND $e_1$ \ldots\ $e_n$
}\\
		If the {\tt AND} expression has more than one argument, and one of
its arguments reduces to {\tt TRUE}, the {\tt TRUE} is dropped from the
expression.\\[.10in]
	{\tt AND $e_1$ \ldots\ FALSE \ldots\ $e_n$ \reds\  FALSE}\\
	  If any of the arguments reduces to {\tt FALSE}, then the entire {\tt AND}
expression reduces to {\tt FALSE}.
\end{quote}
The arguments are reduced in order from left to right.

	Implied by these rules, but not explicitly stated, is that some of the
arguments might not reduce to a boolean value; these non-boolean
arguments remain in place, and reduction of the {\tt AND} continues.	
The reasoning behind this behavior has to do with strong head normalization: 
it is desirable to squeeze out as much information (i.e., perform as
many reductions) as possible from an expression.  Here are a few examples:
\begin{quote}
\begin{tt}
(and ($<$ 3 4) (odd? 5)) \reds\\ 
(and true (odd? 5)) \reds\\ 
(and (odd? 5)) \reds\\
(and true) \reds\\
true\\[.10in]
(and ($>$ 3 4) (odd? 5)) \reds\\  
(and false (odd? 5)) \reds\\
false\\[.10 in]
(and ($<$ 3 4) (even? 5)) \reds\\ 
(and true (even? 5)) \reds\\ 
(and (even? 5)) \reds\\
(and false) \reds\\
false\\[.10in]
(and ($>$ x 4) (odd? 5)) \reds\\
(and ($>$ x 4) true) \reds\\
(and ($>$ x 4))
\end{tt}
\end{quote}


\subsection{\tt OR}

	The semantics of the {\tt OR} primitive are similar to those of the
{\tt AND}, and can be expressed by the
following rules:
\begin{quote}
	{\tt OR FALSE\reds\  FALSE}\\
		 If {\tt OR} has a single argument which
reduces to {\tt FALSE}, then the {\tt OR} expression reduces to
{\tt FALSE}\\[.10in]
	{\tt OR $e_1$ \ldots\ FALSE \ldots\ $e_n$ \reds\  OR $e_1$ \ldots\ $e_n$
}\\
		If the {\tt OR} expression has more than one argument, and one of
its arguments reduces to {\tt FALSE}, the {\tt FALSE} is dropped from the
expression.\\[.10in]
	{\tt OR $e_1$ \ldots\ TRUE \ldots\ $e_n$ \reds\  TRUE}\\
	  If any of the arguments reduces to {\tt TRUE}, then the entire {\tt OR}
expression reduces to {\tt TRUE}.
\end{quote}
The comments made about the {\tt AND} primitive also apply to {\tt OR}.

\section{The {\tt IF} Conditional}

\begin{quote}
	{\tt (IF {\sl cond} {\sl true-exp} {\sl false-exp})}
\end{quote}

	There is only one conditional operator in THOR, the {\tt IF}.
The {\tt IF} is strict only in its first argument, the conditional
expression {\sl cond}.  If {\sl cond} reduces to {\tt TRUE}, then the 
{\tt IF} expression reduces to {\sl true-exp};  if {\sl cond} reduces to
{\tt FALSE}, then the {\tt IF} expression reduces to {\sl false-exp}. 
If {\sl cond} does not reduce to either {\tt TRUE} or {\tt FALSE}, then
the {\tt IF} expression is left as is, and neither {\sl true-exp} nor
{\sl false-exp} is reduced.


\section{Recursion}

	There are three forms of recursion in THOR:  explicit single
recursion using the {\tt REC} operator; implicit single recursion
accomplished through symbols that have associated definitions; and
explicit multiple recursion using the {\tt LETREC} abstraction
mechanism.

\subsection{The {\tt REC} Operator}

	The {\tt REC} operator is used to explicitly cause recursion.  It is a
strange and wondrous thing, with the following semantics:
\begin{quote}\begin{tt}
	(REC {\sl f}) \reds\ ({\sl f} (REC {\sl f}))
\end{tt}\end{quote}
where {\sl f} is an arbitrary expression.  An expression for computing
the factorial of an integer is presented as an example of recursion using
the {\tt REC} operator:
\begin{center}\begin{tt}
	(rec (lambda (fact n) (if (= n 0) 1 (* n (fact (1- n))))))
\end{tt}\end{center}
Note that {\tt fact} is a variable, which gets bound to the
entire {\tt REC} expression during reduction.

	One must be careful when
using the {\tt REC} operator (and any of the other forms of recursion) not
to create an expression which has no normal form, or the reduction
sequence will not terminate.

\subsection{Recursion via Symbol Definitions}

	Recursion can be obtained implicitly through the reduction of symbols
that have been associated with expressions by the use of a defining form
(see Section~\ref{sec-defining}).  This is the form of recursion most
people are familiar with.  To contrast this method with using the {\tt
REC} operator, here is another expression for computing factorials:
\begin{center}\begin{tt}
	(define fact  (lambda (n) (if (= n 0) 1 (* n (fact (1- n))))))
\end{tt}\end{center}
In this expression, {\tt fact} is not a variable, but a symbolic
constant (that just happens to be on both sides of the definition).  
When the symbol {\tt
fact} is encountered as an operator during reduction, it is replaced by
the expression 
\begin{center}
	{\tt (lambda (n) (if (= n 0) 1 (* n (fact (1- n)))))}
\end{center}


\pagebreak
\subsection{Mutual Recursion with {\tt LETREC}}

\begin{quote}
	{\tt (LETREC ((${\mbox{\sl var}}_1$ ${\mbox{\sl exp}}_1$) \ldots) 
				{\sl body})}
\end{quote}

	The {\tt LETREC} abstraction mechanism provides a way to create local
mutually recursive expressions.  Each of the binding pairs are bound in
parallel in a context that has been specially extended to include all of
the binding pairs, so that any ${\mbox{\sl var}}_i$ can be referenced
inside any ${\mbox{\sl exp}}_j$.  The {\sl body} is then reduced in this
new context. There is no reduction charge for creating the recursion
context, but there is a charge of one reduction each time a 
${\mbox{\sl var}}_i$ references this context.

	As a highly contrived example of mutual recursion, the following
expression determines if an integer is even:
\begin{tt}
\begin{tabbing}
12\=1234\=12345\=\kill
   (lambda (x)\\
\>        (letrec ((even? (lambda (n) (if (= n 0) TRUE (odd? (1- n)))))\\
\>\>\>               (odd? (lambda (n) (if (= n 0) FALSE (even? (1-
n))))))\\
\>\>			(even? x)))
\end{tabbing}
\end{tt}

	
	When in step mode, the reduction semantics of the {\tt LETREC} 
are more complex.  The behavior of {\tt LETREC} must take into account
THOR's requirement that the user be able to control the number of
reductions performed on an expression.
The special context of local recursive definitions that is created by
a letrec may need to exist for however long the letrec's
body exists, which may be many reductions. But the user may wish to stop
reduction before the body is completely reduced.
Therefore, some sort of 
mechanism must exist for preserving the recursive context across
invocations of the reducer.  This is done by reconstructing the letrec
binding expression everywhere its context will be needed again.  These
places are easy to find, because they are simply everywhere one of the 
${\mbox{\sl var}}_i$ exist in the body.  

	An example of this behavior is in order.  The expression
\begin{tt}
\begin{tabbing}
12345\=111\=111111\=\kill
\>(letrec ((x (cons 1 y))\\
\>\>\>	  (y (cons 2 x)))\\
\>\>	x)
\end{tabbing}
\end{tt}
describes an infinite list of alternating 1's and 2's.  Assume that
HORSE is in step mode and that the force switch is turned on, so that lists
evaluate eagerly.  Allowing a single reduction lets the {\tt LETREC}
create the recursion context and have the body, which is the recursion 
variable {\tt x}, be expanded by its value, {\tt (cons 1 y)}.  But now
there are no more reductions allowed, and all of the recursion variables
in the body must be `wrapped' in a reconstructed letrec.  (There happens
to be just one recursion variable in the body, the {\tt y}.)  So the
expression is now:
\pagebreak
\begin{tt}
\begin{tabbing}
12345\=22222\=111\=111111\=\kill
\>(cons 1\\
\>\>(letrec ((x (cons 1 y))\\
\>\>\>\>	  (y (cons 2 x)))\\
\>\>\>	y))
\end{tabbing}
\end{tt}
Allowing another reduction causes the {\tt CONS} to be reduced:
\begin{center}\begin{tt}
[1 $\parallel$ (letrec ((x (cons 1 y)) (y (cons 2 x))) y)]
\end{tt}\end{center}
The cdr of this list is now of a form similar to that of the original
expression, and its reduction sequence unfolds in a similar way.  Another
reduction gives
\begin{center}\begin{tt}
[ 1 $\parallel$ (cons 2 (letrec ((x (cons 1 y)) (y (cons 2 x))) x))]
\end{tt}\end{center}
The reduction sequence can continue on in this fashion, ad infinitum.



\section{The {\tt EQUAL?} Primitive}

	In the lambda calculus, two expressions are $\alpha$-equivalent if
they are syntactically the same except for their bound variable names. 
This is a way of saying that the syntactic structure of the expressions is
identical.  The {\tt EQUAL?} predicate tests for this kind of
structural equality.

	{\tt EQUAL?} is a strict binary operator that works by decomposing
its arguments into their atomic constituents and then checking these
for syntactic
equality. If all the parts are equal, the result is {\tt TRUE};  if the
parts are not equal, the result is {\tt FALSE}; and if some of the parts
are unbound variables, then the result will be a new expression that is
a conjunction of equality tests on the unbound variables.  The process
of decomposing expressions and checking their constituent 
subexpressions for $\alpha$-equality may take many reductions, so some
strange looking expressions may be generated as intermediate forms. 
Also, if any of the constituents are lazy compound data objects, the
elements of the objects will be reduced before they are compared.

	Here are some examples:
\begin{tt}\begin{tabbing}
12345\=12345\=\kill
\>	(equal? 1 1) \red\  true\\
\>	(equal? 1 "hello") \red\  false\\
\>	(equal? (a b) (a b)) \red\ true\\
\>	(equal? (a b c) (a b)) \red\ false\\
\>	(lambda (x) (equal? (a 1) (x 1))) \red\\
\>\>		 (lambda (x) (equal? a x))\\
\>	(lambda (x y) (equal? (x y) (a b))) \red\\
\>\>	 (lambda (x y) (and (equal? x a) (equal? y b))\\
\>	(equal? (lambda (a) a) (lambda (b) b)) \red\ true\\
\>	(equal? (lambda (x) 1) (lambda (x) x)) \red\ false\\
\>	(equal? [1 2 3] [1 2 3]) \red\ true\\
\>	(equal? [1 2] [1 2 3]) \red\ false\\
\>	(lambda (x) (equal? [1 x 3] [1 2 3])) \red\\
\>\>	 (lambda (x) (equal? x 2))\\
\>	(equal? \{node 2 (+ 1 4)\} \{node 2 5\}) \red\ true\\
\>	(equal? \{node 2 sym\} \{tree 3 [1 2]\}) \red\ false
\end{tabbing}\end{tt}

	
